\chapter{Case Study: Building an FFI and Native Extensions}
\label{ch:ffi-case-study}

The case study connects one native contract to two managed runtimes and exercises the ABI independently of either runtime.
Its purpose is to make the obligations in the preceding chapters executable: exact integer conversion, explicit byte lengths, matching allocation and release, protected callbacks, and address-stable call metadata.
The implementation lives in \code{example-code/ffi-labs/}.
It uses C11, libffi, Lua 5.4, and the CPython C API, with Python models for selected layout and lifetime protocols.

\section{Keep the supported surface explicit}
\label{sec:ch76-scope}
The reusable scalar call-plan API supports four kinds and at most eight arguments.
Its targets must have trusted matching signatures, use the host default ABI, and remain loaded throughout execution.
It does not infer types from symbols or expose arbitrary user-selected addresses through Lua or Python.
The public managed operations are checked addition, copying a byte buffer into native ownership, copying it back, closing it, and visiting integer indices with a synchronous callback.

Table~\ref{tab:ch76-files} maps these responsibilities to files.
The native tests additionally use libffi directly for two aggregate types, promoted variadic arguments, and one executable closure.
These fixtures verify selected ABI paths without expanding the scalar adapter's public promise.

\begin{table}[htbp]
\caption{The core contract is shared while each runtime owns its adaptation rules.}
\label{tab:ch76-files}
\centering
\begin{tabular}{ll}
\toprule
File & Responsibility \\
\midrule
\code{bookffi.h}, \code{bookffi.c} & C core, scalar plans, versioned API table \\
\code{python\_module.c} & Limited API module and capsule ownership \\
\code{lua\_module.c} & Lua functions, userdata, protected callbacks \\
\code{lua\_host.c} & embedded Lua host with protected evaluation \\
\code{native\_test.c} & direct-call/FFI comparisons and native lifetimes \\
\code{models.py} & restricted layout and sequential protocol models \\
\bottomrule
\end{tabular}
\end{table}

The tested target is Linux x86-64 with the ordinary System V ABI, 32-bit C \code{int}, 64-bit pointers, and the installed libffi.
The Lua configuration uses 64-bit integers.
The Python extension selects a 3.11 Limited API floor and is tested on the installed GIL-enabled CPython; no cross-version binary matrix is inferred from that header setting.
Detailed dependency versions and test commands appear in the lab README and the expansion verification record.

\section{Build a core with failure-atomic outputs}
\label{sec:ch76-core}
The checked add function computes its intermediate in signed 64-bit arithmetic.
The largest possible magnitude from adding two signed 32-bit values fits that type.
Only after checking the result against the signed 32-bit bounds does it write the output.
Listing~\ref{lst:ch76-add} is included directly from the tested source.

\begin{listing}[htbp]
\caption{The native core publishes an integer only after validating the result range.}
\label{lst:ch76-add}
\inputminted[firstline=8,lastline=16]{c}{example-code/ffi-labs/bookffi.c}
\end{listing}

The buffer constructor checks the size limit and allocation-size addition before allocating.
It copies only a nonzero count and permits a null data pointer for the zero-length case.
The pointer returned by \code{bb\_data} is borrowed from a live buffer and becomes invalid after release.
The managed wrappers never expose that pointer directly; they return a fresh language-owned byte value.

Figure~\ref{fig:ch76-buffer} follows the Python constructor.
The input export is released after the copy even if later capsule creation fails.
If capsule creation fails, the payload and owner are freed; if it succeeds, the capsule destructor owns those obligations.
The tests resize the original bytearray afterward to confirm the export is no longer outstanding.

\begin{figure}[htbp]
\centering
\begin{tikzpicture}[node distance=11mm]
\node[basic block] (export) {acquire contiguous input export};
\node[basic block,below=of export] (copy) {allocate owner and copy bytes};
\node[basic block,below=of copy] (release) {release temporary input export};
\node[basic block,below=of release] (publish) {publish capsule or free native ownership};
\draw[cfg edge] (export)--(copy); \draw[cfg edge] (copy)--(release);
\draw[cfg edge] (release)--(publish);
\end{tikzpicture}
\caption{The temporary input borrow ends independently of the published native object's lifetime.}
\label{fig:ch76-buffer}
\end{figure}

The core functions assume valid native pointers and nonoverlapping ranges for their copies.
They reject nulls and inconsistent counts where specified, but cannot prove arbitrary pointer validity from numerical values.
In particular, the caller must not read or free a buffer after release, nor provide an undersized output region to a successful read.
Opaque managed owners enforce a narrower accessible state space than unrestricted native callers have.

\section{Prepare and invoke one known signature}
\label{sec:ch76-plan}
Checked addition has three arguments: two signed 32-bit values and a pointer to the output integer.
Its status is an ordinary 32-bit C integer on the tested target.
The wrapper creates typed argument slots, prepares a call interface, and invokes the known core function through libffi.
Algorithm~\ref{alg:ch76-add} follows this complete path from a managed value pair to a managed result.

\begin{algorithm}[htbp]
\caption{Execute checked managed addition through the scalar foreign-call bridge.}
\label{alg:ch76-add}
\begin{pseudocode}
\Require two language values; supported exact-integer conversion policy
\Ensure managed sum or a language error with no uninitialized result read
\State check both values' kinds and signed 32-bit ranges
\State prepare address-stable plan for two int32 arguments and one pointer
\State create typed slots and local result storage
\State invoke the known native checked-add target through libffi
\If{bridge failure or native status is not success}
  \State translate status; do not read the arithmetic output
\Else
  \State convert initialized arithmetic output to a language integer
\EndIf
\end{pseudocode}
\end{algorithm}

For this fixed signature, adapter bookkeeping is constant apart from language integer conversion.
The example prepares a fresh plan for each call to keep its ownership local and avoid a shared mutable cache.
A production cache would key on the complete ABI and signature and retain immutable, address-stable metadata until all users finish.
Caching by argument count alone would conflate, for example, two doubles with two pointers.

\section{Bind ownership and errors separately in each runtime}
\label{sec:ch76-bindings}
The Python constructor obtains a contiguous buffer export, copies it, releases the export, and constructs a named capsule retaining the native owner.
The close function clears only the payload; the capsule continues owning the small owner record until destruction.
Every subsequent byte-read checks the payload for the closed state.
The module counter lives in per-module state allocated by multi-phase initialization~\cite{python-2026-capi}.

The Lua constructor accepts a string without numeric coercion, creates full userdata with a null payload, installs its cleanup metatable, and then creates the native copy.
The same close method implements explicit close, garbage-collection cleanup, and Lua's to-be-closed scope handling.
The input string remains rooted on the stack throughout construction~\cite{ierusalimschy-2026-lua}.
The embedded host registers this same module implementation and evaluates its test script through an outer protected call.

Figure~\ref{fig:ch76-errors} shows the error adapters around the shared visitor.
For Python, a failed callback leaves the exception indicator set and the outer wrapper returns null.
For Lua, the protected call retains the original error object on the stack and the outer wrapper re-raises it after native return.
Both avoid publishing a partial visitor sum.

\begin{figure}[htbp]
\centering
\begin{tikzpicture}
\node[basic block] (native) at (0,0) {native visitor returns status};
\node[basic block,text width=4cm] (python) at (-3,-1.6) {Python: null plus exception};
\node[basic block,text width=4cm] (lua) at (3,-1.6) {Lua: re-raise retained object};
\draw[cfg edge] (native)--(python); \draw[cfg edge] (native)--(lua);
\end{tikzpicture}
\caption{A common native status maps back to each runtime's own error mechanism.}
\label{fig:ch76-errors}
\end{figure}

\section{Test independent failure dimensions}
\label{sec:ch76-tests}
A round trip can hide matching bugs on both sides, so the native suite compares selected libffi calls against compiler-generated direct calls.
It also compares structure offsets, sizes, and alignment against C's own operators.
Boundary tests use a mathematical widened sum to distinguish valid arithmetic from overflow, and verify that failure preserves sentinels.
Algorithm~\ref{alg:ch76-differential} describes the native differential method.

\begin{algorithm}[htbp]
\caption{Check a foreign-call path against a compiler-generated native oracle.}
\label{alg:ch76-differential}
\begin{pseudocode}
\Require trusted declaration, native function, ABI descriptors, boundary-value cases
\Ensure compared paths agree on status and defined output fields
\ForAll{input cases}
  \State initialize independent native and FFI output storage with sentinels
  \State call the declared function directly
  \State call the same function through libffi with independently built slots
  \State compare statuses and every semantically defined result field
  \State verify unchanged outputs where the failure contract requires them
\EndFor
\State compare descriptor layout with native size, alignment, and offsets
\end{pseudocode}
\end{algorithm}

For $T$ tests and $A$ arguments, setup and comparison cost $O(TA)$ for fixed-size values, excluding the calls themselves.
The method does not prove libffi correct for every signature.
Direct and foreign paths share the same native function, so independent arithmetic expectations and lifetime tests remain necessary.
Structure comparison uses fields, not padding bytes.

Table~\ref{tab:ch76-tests} separates executed native mechanisms from bounded models.
The registration model enumerates all $4^7=16{,}384$ event words of length seven, checking each legal prefix and rejecting illegal transitions.
It does not test actual managed callbacks from foreign threads.

\begin{table}[htbp]
\caption{The test suite checks ABI, runtime, and protocol obligations through different oracles.}
\label{tab:ch76-tests}
\centering
\begin{tabular}{ll}
\toprule
Check & Evidence \\
\midrule
scalar ABI & 49 boundary pairs; mixed banks; eight arguments \\
aggregate ABI & Pair and Triple fieldwise direct-call comparisons \\
variadic ABI & promoted doubles; unpromoted float rejection \\
closure entry & 101 calls with retained bias context \\
native threads & four workers, 1,000 independent calls each \\
runtime ownership & copies, closed-state rejection, repeated close, collection \\
runtime callbacks & nesting, conversion failure, original error identity \\
protocol model & all seven-event registration words; moving handles \\
\bottomrule
\end{tabular}
\end{table}

The build matrix uses GCC and Clang, each with normal and address/undefined-behavior-sanitized native executables.
The embedded Lua host runs in both configurations; dynamically loaded Lua and Python extensions run in ordinary interpreter processes without injecting sanitizer runtimes.
This distinction prevents a sanitized-host result from being misreported as a sanitized CPython-extension result.
The suite also checks plugin negotiation and closes its loader handle after the last possible use.

\section{Extend only with a new contract and its checks}
\label{sec:ch76-extension}
Adding asynchronous callbacks requires retained roots, thread attachment or queued delivery, unregister acknowledgement, and a tested reclamation protocol.
Adding zero-copy buffers requires an export owner, mutation rules, and coordination between outstanding views and close.
Adding arbitrary aggregate signatures requires a type-description representation and lifetime discipline beyond the four-kind scalar union.
Each is a meaningful extension, but none follows merely from adding a function name to the module table.

The tests do not inject every allocator failure, test every operating system, or validate a cross-version Stable ABI distribution.
Those are concrete next verification dimensions rather than claims established by the current matrix.
The implemented result is a small interoperable system with explicit limits and independently checked examples of the mechanisms that broader FFIs need.

\section{Exercises}
\label{sec:ch76-exercises}
\begin{enumerate}
\item\label{ex:ch76-plan} Explain why copying a prepared \code{FxPlan} with a structure assignment can invalidate its internal ownership assumptions.
\item\label{ex:ch76-error} For each of input $2^{80}$, addition overflow, and callback exception, identify the first rejecting component and the error returned to Python.
\item\label{ex:ch76-cost} A buffer of $N$ bytes is constructed and immediately read back.
Count the payload copies and their asymptotic time and storage costs.
\item\label{ex:ch76-cache} Specify a safe cache key and metadata lifetime rule for reusing prepared scalar call plans across invocations.
Explain why argument count is insufficient.
\end{enumerate}

\section{Further reading}
\label{sec:ch76-reading}
The libffi manual supplies the ABI adapter contracts used by the implementation~\cite{libffi-2026-manual}.
The Lua and Python manuals define the runtime mechanisms that complete the ownership and error boundary~\cite{ierusalimschy-2026-lua,python-2026-capi}.
The lab README records reproducible commands, and its solution file works every exercise from the eight FFI chapters.
